%==============================================================
%  preamble.tex -- packages, colours, helper macros.
%==============================================================
\usepackage[T1]{fontenc}
\usepackage[utf8]{inputenc}
\usepackage[a4paper,margin=1in,headheight=26pt]{geometry}
\usepackage{setspace}
\onehalfspacing
\IfFileExists{lmodern.sty}{\usepackage{lmodern}}{}
\usepackage[expansion=false]{microtype}

% ---- graphics, colours, tables ---------------------------------
\usepackage[dvipsnames,table]{xcolor}
\usepackage{graphicx}
\graphicspath{{images/}{images/diagrams/}{images/photos/}{images/ss/}}
\usepackage{float}
\usepackage{booktabs,array,tabularx,longtable,multirow,makecell}
\usepackage{caption,subcaption}
\captionsetup{font=small,labelfont=bf}
\usepackage{enumitem}
\setlist{itemsep=2pt,topsep=3pt}
\usepackage{amsmath,amssymb}
\usepackage{pifont}
\newcommand{\cmark}{\ding{51}}
\usepackage{pdflscape}

% ---- CUET colour palette ---------------------------------------
\definecolor{cuetblue}{RGB}{38,36,115}
\definecolor{cuetgold}{RGB}{218,165,32}
\definecolor{cuetgreen}{RGB}{34,139,34}

% ---- diagrams / Gantt ------------------------------------------
\usepackage{tikz}
\usetikzlibrary{shapes,arrows.meta,positioning,fit,calc}
\usepackage{pgfgantt}

% ---- code listings ---------------------------------------------
\usepackage{listings}
\lstdefinestyle{cuetcode}{
  basicstyle=\ttfamily\footnotesize,
  keywordstyle=\color{cuetblue}\bfseries,
  commentstyle=\color{cuetgreen}\itshape,
  stringstyle=\color{red!60!black},
  numbers=left, numberstyle=\tiny\color{gray}, numbersep=6pt,
  frame=single, rulecolor=\color{gray!60},
  breaklines=true, breakatwhitespace=false,
  showstringspaces=false, tabsize=4, captionpos=b,
  backgroundcolor=\color{gray!5}}
\lstset{style=cuetcode}
\renewcommand{\lstlistingname}{Listing}

% ---- algorithms (pseudocode for LLD) ---------------------------
\usepackage{algorithm}
\usepackage{algpseudocode}

% ---- headers / footers -----------------------------------------
\usepackage{fancyhdr}
\pagestyle{fancy}
\fancyhf{}
\fancyhead[L]{\small\textcolor{cuetblue}{Infra Monitor -- Project Report}}
\fancyhead[R]{\small\textcolor{cuetblue}{CSE-356}}
\fancyfoot[C]{\thepage}
\renewcommand{\headrulewidth}{0.4pt}
\renewcommand{\headrule}{\hbox to\headwidth{\color{cuetgold}\leaders\hrule height \headrulewidth\hfill}}
\fancypagestyle{plain}{\fancyhf{}\fancyfoot[C]{\thepage}\renewcommand{\headrulewidth}{0pt}}

% ---- chapter / section title colours ---------------------------
\usepackage{titlesec}
\titleformat{\chapter}[display]
  {\normalfont\huge\bfseries}
  {\chaptertitlename\ \thechapter}{10pt}{\Huge}
\titlespacing*{\chapter}{0pt}{-20pt}{25pt}
\titleformat{\section}{\normalfont\Large\bfseries\raggedright}{\thesection}{1em}{}
\titleformat{\subsection}{\normalfont\large\bfseries}{\thesubsection}{1em}{}

% ---- figure helper: prints a placeholder if the image is missing
\newcommand{\figplaceholder}[4][0.85\textwidth]{%
  \begin{figure}[H]\centering
    \IfFileExists{images/diagrams/#2}{\includegraphics[width=#1]{#2}}%
    {\IfFileExists{images/#2}{\includegraphics[width=#1]{#2}}%
    {\fbox{\parbox{#1}{\centering\vspace{2cm}\textcolor{red}{[Figure: \texttt{\detokenize{#2}}]}\vspace{2cm}}}}}
    \caption{#3}\label{#4}
  \end{figure}}

% ---- team member macros ----------------------------------------
\makeatletter
\newcommand{\coverlist}{}
\newcommand{\tablerows}{}
\newcommand{\memberphoto}[1]{%
  \IfFileExists{#1}{\includegraphics[width=2cm,height=2.4cm,keepaspectratio]{#1}}%
  {\includegraphics[width=2cm,height=2.4cm,keepaspectratio]{images/photos/placeholder.png}}}
\newcommand{\member}[7]{%
  \g@addto@macro\coverlist{#1\\ ID: #2\\[1.4ex]}%
  \g@addto@macro\tablerows{\memberphoto{#7} & #1 & #2 & {\scriptsize\url{#3}} & #4 & #5 & #6 \\ \hline}}
\makeatother

% ---- hyperlinks (load LAST) ------------------------------------
\usepackage[hidelinks]{hyperref}
\hypersetup{pdftitle={Infra Monitor -- Software Engineering Final Report}}
\usepackage[nameinlink]{cleveref}

% ---- handy macros ----------------------------------------------
\newcommand{\todo}[1]{\textcolor{red}{\textbf{[#1]}}}
\newcommand{\req}[1]{\texttt{#1}}
\newcommand{\assumed}[1]{\textcolor{orange}{\textbf{[ASSUMED: #1]}}}
\newcommand{\tofill}[1]{\textcolor{red}{\textbf{[TO FILL: #1]}}}

